\section{Results}



\subsection{Experimental Overview}



The final benchmark evaluated four classification models—XGBoost, LightGBM, CatBoost, and TabPFN—across three binary tabular classification datasets: Adult, Bank Marketing, and Credit-G. Five random seeds were evaluated for each training-size configuration.



Adult and Bank Marketing were evaluated at 21 training sizes ranging from 10 to 32,000 observations. Credit-G was evaluated at 10 training sizes ranging from 10 to 750 observations because of its smaller available training partition. Across all datasets, models, training sizes, and random seeds, the benchmark produced 1,040 model-level observations.



Predictive performance was evaluated using Accuracy, Precision, Recall, F1-score, and ROC-AUC. Training time and prediction time were also recorded.



The results are presented from four perspectives: overall predictive performance, paired statistical comparisons, performance across training-data regimes, and computational performance.



\subsection{Overall Predictive Performance}



The aggregate results show that no single model dominates across all datasets and evaluation metrics. The strongest model varies according to both the dataset and the metric.



\subsubsection{Adult}



For Adult, CatBoost achieves the highest mean Accuracy at \textbf{0.8409}, Precision at \textbf{0.7330}, F1-score at \textbf{0.5936}, and ROC-AUC at \textbf{0.8787}.



XGBoost achieves the highest mean Recall at \textbf{0.5334}.



TabPFN achieves mean Accuracy of \textbf{0.8350}, Precision of \textbf{0.7031}, Recall of \textbf{0.5238}, F1-score of \textbf{0.5838}, and ROC-AUC of \textbf{0.8744}.



Thus, TabPFN performs competitively on Adult but does not achieve the strongest aggregate performance. CatBoost remains the strongest overall model on this dataset across most of the evaluated metrics.



The relatively small difference between TabPFN and CatBoost in ROC-AUC is notable: TabPFN achieves 0.8744 compared with 0.8787 for CatBoost. However, the aggregate results do not show an overall TabPFN advantage on Adult.



\subsubsection{Bank Marketing}



Bank Marketing provides the strongest overall results for TabPFN.



TabPFN achieves the highest mean Accuracy (\textbf{0.8980}), Recall (\textbf{0.3201}), F1-score (\textbf{0.3860}), and ROC-AUC (\textbf{0.8716}).



CatBoost achieves the highest Precision (\textbf{0.6022}).



The difference in ROC-AUC is particularly notable. TabPFN achieves a mean ROC-AUC of 0.8716 compared with 0.8581 for CatBoost, 0.8145 for XGBoost, and 0.8115 for LightGBM.



Similarly, TabPFN achieves a mean F1-score of 0.3860 compared with 0.3464 for CatBoost, 0.3645 for XGBoost, and 0.3479 for LightGBM.



Therefore, Bank Marketing provides the clearest evidence in the benchmark that TabPFN can outperform established gradient-boosted tree models across multiple predictive dimensions.



\subsubsection{Credit-G}



Credit-G produces a different pattern.



CatBoost achieves the highest mean Accuracy (\textbf{0.7212}), Recall (\textbf{0.8957}), F1-score (\textbf{0.8182}), and ROC-AUC (\textbf{0.7213}).



XGBoost achieves the highest Precision (\textbf{0.7621}).



TabPFN achieves Accuracy of \textbf{0.7145}, Precision of \textbf{0.7550}, Recall of \textbf{0.8859}, F1-score of \textbf{0.8126}, and ROC-AUC of \textbf{0.7146}.



The differences between TabPFN and CatBoost are relatively small for several metrics. For example, the F1-score difference is approximately 0.0055 in favour of CatBoost, while the ROC-AUC difference is approximately 0.0067.



Consequently, Credit-G does not provide evidence of a clear overall TabPFN advantage. Instead, it demonstrates that a conventional gradient-boosted tree model can remain highly competitive on smaller tabular datasets.



\subsection{Statistical Comparison}



To determine whether observed differences between TabPFN and the tree-based baselines were consistent across repeated experimental configurations, paired comparisons were performed using observations matched by dataset, training size, and random seed.



There were 45 TabPFN-versus-tree comparisons in total: three tree-based baselines across five evaluation metrics and three datasets.



Of these 45 comparisons, \textbf{34 were statistically significant at the unadjusted 0.05 level} according to the Wilcoxon signed-rank test. After Holm correction across all 45 comparisons, \textbf{31 remain statistically significant at the 0.05 level}. The paired-results table reports both unadjusted and Holm-adjusted p-values and identifies the Holm-corrected significance status.



Statistical significance alone does not indicate which model performed better. The direction and magnitude of the paired differences were therefore also examined.



\subsubsection{Adult}



TabPFN significantly outperforms XGBoost and LightGBM in Accuracy, Precision, F1-score, and ROC-AUC, although its Recall is significantly lower than XGBoost.



For ROC-AUC, TabPFN exceeds XGBoost by \textbf{0.0280} and LightGBM by \textbf{0.0521}. Both differences are statistically significant and have large rank-biserial effect sizes.



However, TabPFN performs significantly worse than CatBoost in Accuracy, Precision, F1-score, and ROC-AUC. The ROC-AUC difference is \textbf{-0.0043} in favour of CatBoost.



The statistical results therefore confirm the aggregate results: TabPFN is competitive with XGBoost and LightGBM on Adult, but CatBoost remains a particularly strong baseline.



\subsubsection{Bank Marketing}



Bank Marketing provides the strongest statistical evidence in favour of TabPFN.



TabPFN significantly outperforms XGBoost and LightGBM across all five metrics.



The ROC-AUC difference is \textbf{0.0571} relative to XGBoost and \textbf{0.0601} relative to LightGBM. Both differences are statistically significant and have very large rank-biserial effect sizes.



TabPFN also significantly outperforms CatBoost in Accuracy, Recall, F1-score, and ROC-AUC.



For ROC-AUC, TabPFN exceeds CatBoost by \textbf{0.0135}. For F1-score, the difference is \textbf{0.0396}.



The only non-significant TabPFN-CatBoost comparison on Bank Marketing is Precision. This is consistent with CatBoost achieving the highest aggregate Precision on the dataset.



\subsubsection{Credit-G}



Credit-G produces a more mixed result.



TabPFN significantly outperforms XGBoost in Accuracy, Recall, F1-score, and ROC-AUC. The corresponding differences are 0.0170, 0.0539, 0.0195, and 0.0403 respectively.



TabPFN also significantly outperforms LightGBM in Accuracy, F1-score, and ROC-AUC.



In contrast, none of the five TabPFN-CatBoost comparisons reaches statistical significance.



This is consistent with the aggregate results, where CatBoost achieves the strongest mean performance across most metrics.



Overall, the statistical analysis demonstrates that the differences between TabPFN and tree-based models are systematic in many experimental configurations, but the direction of those differences depends strongly on the dataset and baseline model.



\subsection{Performance Across Training-Data Regimes}



The learning-curve analysis examines how model performance changes as the amount of training data increases.



For interpretation, the analysis groups the available training sizes into small, medium, and large regimes. Adult and Bank Marketing contain all three regimes, whereas Credit-G contains only small and medium regimes because its training partition does not support the larger sample sizes.



\subsubsection{Adult}



In the small-data regime, TabPFN performs better than XGBoost and LightGBM on Accuracy, Precision, F1-score, and ROC-AUC. For example, its mean ROC-AUC is \textbf{0.7993}, compared with \textbf{0.7304} for XGBoost and \textbf{0.6419} for LightGBM.



However, CatBoost remains competitive, with a mean small-regime ROC-AUC of \textbf{0.8015}.



In the medium regime, TabPFN again exceeds XGBoost and LightGBM on most metrics. Its mean ROC-AUC is \textbf{0.8900}, compared with \textbf{0.8629} for XGBoost and \textbf{0.8589} for LightGBM.



At large training sizes, the differences become smaller and the tree-based models become increasingly competitive. CatBoost achieves higher mean Accuracy, Recall, F1-score, and ROC-AUC than TabPFN in this regime. For ROC-AUC, CatBoost reaches \textbf{0.9182} compared with \textbf{0.9116} for TabPFN.



The Adult learning-regime results therefore show a clear change in relative performance as training data increase: TabPFN is competitive or stronger against XGBoost and LightGBM in smaller regimes, while CatBoost becomes stronger in the large-data regime.



\subsubsection{Bank Marketing}



Bank Marketing shows the most consistent TabPFN advantage across the three training regimes.



In the small regime, TabPFN achieves a mean ROC-AUC of \textbf{0.7520}, compared with \textbf{0.6437} for XGBoost and \textbf{0.6195} for LightGBM.



In the medium regime, TabPFN achieves a mean ROC-AUC of \textbf{0.8907}, compared with \textbf{0.8364} for XGBoost and \textbf{0.8388} for LightGBM.



In the large regime, the advantage remains visible. TabPFN reaches a mean ROC-AUC of \textbf{0.9339}, compared with \textbf{0.9060} for XGBoost, \textbf{0.9131} for LightGBM, and \textbf{0.9225} for CatBoost.



A similar pattern is observed for F1-score. In the large regime, TabPFN achieves \textbf{0.5521}, compared with \textbf{0.4963} for XGBoost, \textbf{0.4957} for LightGBM, and \textbf{0.5077} for CatBoost.



These results indicate that the TabPFN advantage on Bank Marketing is not restricted to the smallest training-data regime. It persists across the evaluated training-size ranges.



\subsubsection{Credit-G}



Credit-G contains only small and medium training regimes.



In the small regime, TabPFN achieves a mean ROC-AUC of \textbf{0.6767}, compared with \textbf{0.6376} for XGBoost and \textbf{0.5814} for LightGBM. CatBoost achieves \textbf{0.6922}, making it the strongest model among the four for this metric in this regime.



In the medium regime, TabPFN reaches a mean ROC-AUC of \textbf{0.7714}, compared with \textbf{0.7292} for XGBoost, \textbf{0.7273} for LightGBM, and \textbf{0.7651} for CatBoost.



For F1-score in the medium regime, TabPFN achieves \textbf{0.8191}, while CatBoost achieves \textbf{0.8245}.



The Credit-G learning-regime results therefore show that TabPFN can outperform XGBoost and LightGBM in several settings while remaining closely matched with CatBoost.



\subsubsection{Overall Training-Regime Pattern}



Taken together, the learning-regime results do not support a simple relationship in which TabPFN is strongest only when very little training data are available.



Instead, three different patterns are observed:



\begin{itemize}

\item \textbf{Adult:} TabPFN is competitive in smaller regimes but CatBoost becomes stronger at larger training sizes.

\item \textbf{Bank Marketing:} TabPFN maintains an advantage across small, medium, and large regimes.

\item \textbf{Credit-G:} TabPFN remains competitive with the tree models, while CatBoost remains particularly strong.

\end{itemize}



Thus, the effect of training-data size is \textbf{dataset-dependent rather than universal}.



\subsection{Computational Performance}



Training and prediction times were recorded for every model and benchmark configuration.



The aggregate computational results show a clear difference between the conventional tree-based models and TabPFN.



For Adult, mean training time is \textbf{0.174 seconds} for XGBoost, \textbf{0.189 seconds} for LightGBM, \textbf{0.558 seconds} for CatBoost, and \textbf{0.679 seconds} for TabPFN.



For Bank Marketing, the corresponding mean training times are \textbf{0.212 seconds}, \textbf{0.205 seconds}, \textbf{0.625 seconds}, and \textbf{0.685 seconds}.



For Credit-G, mean training times are \textbf{0.065 seconds} for XGBoost, \textbf{0.052 seconds} for LightGBM, \textbf{0.351 seconds} for CatBoost, and \textbf{0.659 seconds} for TabPFN.



Therefore, within this benchmark, TabPFN has the highest mean training time across all three datasets.



The difference is considerably larger for prediction time.



On Adult, mean prediction time is approximately \textbf{0.062 seconds} for XGBoost, \textbf{0.099 seconds} for LightGBM, \textbf{0.049 seconds} for CatBoost, and \textbf{9.036 seconds} for TabPFN.



On Bank Marketing, the corresponding values are \textbf{0.046 seconds}, \textbf{0.091 seconds}, \textbf{0.042 seconds}, and \textbf{8.946 seconds}.



On Credit-G, prediction time is approximately \textbf{0.010 seconds} for XGBoost, \textbf{0.011 seconds} for LightGBM, \textbf{0.011 seconds} for CatBoost, and \textbf{0.667 seconds} for TabPFN.



These results show that TabPFN's predictive performance advantages are accompanied by substantially higher measured prediction times in this experimental implementation.



The computational results therefore provide an important qualification to the predictive results. TabPFN may provide stronger predictive performance on some datasets, but that advantage must be considered alongside its computational cost, particularly during prediction.



\subsection{Summary of Results}



The benchmark does not support the hypothesis that TabPFN universally outperforms gradient-boosted tree algorithms.



Instead, the results show three distinct patterns.



First, \textbf{Bank Marketing provides the clearest evidence in favour of TabPFN}. TabPFN achieves the highest aggregate Accuracy, Recall, F1-score, and ROC-AUC and significantly outperforms the tree-based models on most corresponding paired comparisons. Its advantage is also visible across small, medium, and large training-data regimes.



Second, \textbf{Adult demonstrates that the relative advantage of TabPFN depends on the baseline and training regime}. TabPFN performs strongly against XGBoost and LightGBM, particularly in smaller and medium regimes, but CatBoost achieves stronger aggregate performance and becomes particularly competitive at larger training sizes.



Third, \textbf{Credit-G demonstrates that strong conventional tree-based performance remains possible on smaller tabular datasets}. TabPFN outperforms XGBoost and LightGBM on several metrics, but CatBoost remains closely matched or superior across most aggregate measures.



The statistical analysis further shows that 34 of the 45 TabPFN-versus-tree comparisons are statistically significant at the unadjusted 0.05 level, with 31 remaining significant after Holm correction. However, statistical significance is not uniformly associated with a TabPFN advantage: the direction of the difference varies across datasets and baseline models.



Finally, the computational measurements show a clear trade-off. TabPFN achieves strong predictive performance in several settings but has higher measured training and prediction times than the conventional tree-based models in this benchmark. The prediction-time difference is particularly pronounced on Adult and Bank Marketing.



Overall, the evidence supports the interpretation that \textbf{TabPFN is a complementary approach rather than a universal replacement for gradient-boosted tree algorithms}. Its relative effectiveness depends on the dataset, evaluation metric, training-data regime, and computational requirements.
